
\documentclass[11pt]{article}
\usepackage[margin=1in]{geometry}
\usepackage{amsmath,amssymb,amsthm}
\usepackage{hyperref}
\usepackage{enumitem}

\title{Referee Report and Salvage Map: ``Curvature--Stable Renormalization and the Mass Gap'' Notes}
\date{}

\begin{document}
\maketitle

\section{Clean document set (this is the surviving corpus)}

The project directory contained multiple overlapping LaTeX drafts and computational logs.
The material that survives scrutiny has been reorganized into the following self-contained documents:

\begin{enumerate}[leftmargin=*]
\item \textbf{Document 01:} \emph{Gauge Generator and Physical Projector on a Periodic $4$D Lattice (Finite Volume).}
Definitions of $\mathcal M_L$, the gauge generator, and the \emph{vacuum} orthogonal projector $P_0$.

\item \textbf{Document 02:} \emph{Linearized Wilson Action Hessian at Vacuum and Its Spectrum (Finite Volume).}
A complete derivation of the vacuum Hessian identity
$H_0=(\beta/6)\Delta_1$ and the finite-volume physical gap
$\lambda_{\min}(H_0|\_{\mathrm{phys}})=\frac{\beta}{6}\,4\sin^2(\pi/L)$.
This is the strongest \emph{proved} statement currently supported.

\item \textbf{Document 03:} \emph{Conditional Pipeline: (Gauge Reduction) $\Rightarrow$ (LSI) $\Rightarrow$ (RG Stability) $\Rightarrow$ (OS Mass Gap).}
The intended mass-gap argument as a sequence of \emph{explicit hypotheses} plus labeled external inputs.
No theorem about Yang--Mills mass gap is proved here; this document states exactly what would be required.

\item \textbf{Document 04:} \emph{Quarantined Computational Material (Non-Proofs): Hessian Spectra, Projectors, HOTRG Jacobians.}
A cleaned record of numerical experiments and code sketches, explicitly quarantined from the proof pipeline.
\end{enumerate}

\paragraph{Dependency ordering.}
Document~01 is prerequisite for Document~02.
Document~03 assumes Documents~01--02 for notation and identifies missing lemmas.
Document~04 is independent evidence-only material.

\section{Classification of claims in the original drafts}

The original drafts (notably \texttt{12-5-25 LEMMA2.txt}, \texttt{12-5-25 LEMMA.txt}, \texttt{12-6-25 LEMMA.txt}, \texttt{LEMMA 12-5-25.txt}) mix definitions, conditional theorems, and non sequiturs.
Below is the maximal salvage classification consistent with the corpus.

\subsection*{Proved statements (in the clean docset)}
\begin{enumerate}[leftmargin=*]
\item \textbf{Vacuum Hessian identity and spectrum.}
At $U\equiv I$, the Wilson action Hessian equals $(\beta/6)\Delta_1$ on lattice $1$-forms (per color),
and the strictly physical finite-volume gap equals $\frac{\beta}{6}\,4\sin^2(\pi/L)$.
(See Document~02.)

\item \textbf{Vacuum projector algebra.}
The Moore--Penrose formula
$P_0 = I - G_0(G_0^\dagger G_0)^+G_0^\dagger$
defines an orthogonal projector onto $\ker(G_0^\dagger)$ at vacuum.
(See Document~01.)
\end{enumerate}

\subsection*{Conditionally correct statements (true if explicit hypotheses hold)}
\begin{enumerate}[leftmargin=*]
\item \textbf{``Coercivity $\Rightarrow$ LSI $\Rightarrow$ spectral gap'' on a genuine quotient diffusion.}
This is a standard consequence of Bakry--\'Emery theory \emph{provided} one has:
(i) a gauge reduction to a manifold $Q_L$,
(ii) a symmetric diffusion generator on $Q_L$,
and (iii) a uniform convexity bound for the descended potential.
In the drafts this is asserted for an ad hoc projected generator; in the clean docset it is stated as a conditional pipeline (Document~03).

\item \textbf{RG stability of LSI under Lipschitz pushforward.}
A statement of the form ``LSI constant degrades by at most $L_{\mathrm{RG}}^2$ under an $L_{\mathrm{RG}}$-Lipschitz coarse map''
is conditionally correct as an external input, \emph{provided} the RG step is indeed a Lipschitz pushforward on the relevant metric-measure space.
(Recorded as EI3 in Document~03.)
\end{enumerate}

\subsection*{Conjectural / heuristic statements (not proved, not falsified by a direct contradiction)}
\begin{enumerate}[leftmargin=*]
\item \textbf{Riccati/Hamilton--Jacobi ``convexity preservation'' in the physical sector.}
The drafts state that viscous Hamilton--Jacobi flow preserves the physical convexity bound.
No closed evolution inequality for the projected Hessian is derived in the corpus, and curvature terms are not controlled.
This remains conjectural.

\item \textbf{HOTRG Jacobian control in a right-invariant frame.}
The idea ``use a right-invariant tangent frame so that Hessians push forward as $H'\approx J^\top HJ$'' is a plausible computational architecture.
A theorem requires a careful treatment of (i) manifold second derivatives, (ii) truncation/SVD nonsmoothness, and (iii) reflection positivity/locality preservation.
\end{enumerate}

\subsection*{Incorrect or structurally invalid statements}
The following items are not salvageable as stated.

\begin{enumerate}[leftmargin=*]

\item \textbf{Configuration space mis-specified as $SU(3)^{\Lambda_4}$ or $SU(3)^{\Lambda_L}$.}

Removed / Downgraded because:
the drafts write $\mathcal M_L = SU(3)^{\Lambda_4}$ while also indexing link variables as $U_\mu(x)$ over directions $\mu\in\{0,1,2,3\}$.
This is a category error: the configuration space must be $SU(3)^{E_L}$ with $E_L=\Lambda_L\times\{0,1,2,3\}$ (or an equivalent link-set convention).

\item \textbf{Uniform-in-$L$ vacuum-neighborhood coercivity as stated.}

Removed / Downgraded because:
Lemma ``Gauge--Projected Coercivity'' asserts existence of $C_W>0$ and a \emph{single} neighborhood $\mathcal N$ of vacuum such that
$P_L(U)\nabla^2 S_L(U)P_L(U)\succeq C_W P_L(U)$ for all $L$ and all $U\in\mathcal N$.
Document~02 proves that at the vacuum itself the strictly physical Hessian gap equals $(\beta/6)\,4\sin^2(\pi/L)\sim c\,L^{-2}$,
so no positive lower bound independent of $L$ can hold near vacuum unless the notion of ``physical'' is altered (e.g.\ excluding small momenta) or the theory is modified by an explicit mass term.

\item \textbf{Claim ``Gauge invariance implies $\nabla^2 S(U)\,G(U)\equiv 0$'' (as an identity for all $U$).}

Removed / Downgraded because:
the drafts assert this as the justification for restricting the Hessian to the physical sector.
Gauge invariance implies that the \emph{first} derivative vanishes along gauge orbits; the Hessian annihilation statement does not follow without additional hypotheses and is generally false away from critical points.
Even if the quadratic form vanishes on gauge tangents, the operator identity needs a connection-dependent argument that is not supplied.

\item \textbf{Bakry--\'Emery $\Gamma_2$ computation for the projected generator.}

Removed / Downgraded because:
the drafts define a projected operator of the form $\tilde Lf=\mathrm{div}(P(U)\nabla f)-\langle P(U)\nabla S,P(U)\nabla f\rangle$
and then write $\Gamma_2(f)=\|P\nabla^2 fP\|^2+\langle P\nabla^2 S P\nabla f,\nabla f\rangle$.
This omits the derivatives of the \emph{variable} projector $P(U)$ and does not establish that $\tilde L$ is a symmetric Markov diffusion on any quotient.
Consequently the claimed curvature--dimension inequality is not derived.

\item \textbf{Derivation of spectral gap from LSI via ``$\mathrm{Var}\le \mathrm{Ent}$''.}

Removed / Downgraded because:
the drafts use $\mathrm{Var}(f)\le \mathrm{Ent}(f^2)$ to deduce a Poincar\'e inequality.
This inequality is false in general.  The implication ``LSI $\Rightarrow$ Poincar\'e'' is true, but requires a separate argument (typically a small-perturbation expansion), which is not provided.

\item \textbf{Identification of a ``mass gap'' as $m_{\mathrm{gap}}\ge \sqrt{C_W}$.}

Removed / Downgraded because:
a lower bound on a \emph{Markov generator} spectral gap (even if established) does not directly imply a Hamiltonian mass gap without an Osterwalder--Schrader/transfer-matrix identification plus reflection positivity and locality.
No such justification is supplied in the corpus, and the dimensional relationship $m\sim \sqrt{C_W}$ is not derived.

\item \textbf{Treating numerics as proof of continuum mass gap.}

Removed / Downgraded because:
\texttt{12-6-25 LEMMA.txt} and related drafts explicitly cite numerical Hessian diagonalization / Langevin relaxation / spectroscopy as ``support'' for the convexity hypothesis and then describe the result as establishing a continuum mass gap.
Numerical experiments do not establish the required uniform-in-$L$, uniform-in-$a$ bounds nor the OS reconstruction hypotheses.
All numerical material has been quarantined into Document~04.

\item \textbf{Hallucinated ``volume-independent lower bound $C_W\approx 8/3$'' from $L=2$ micro-tests.}

Removed / Downgraded because:
the project logs also contain computed vacuum gaps $2/3$ at $L=2$ and $1/2$ at $L=3$ in the projected sector, consistent with the analytic $L^{-2}$ scaling.
Thus claims of an $L$-independent lower bound inferred from a single small-$L$ computation are contradicted by other material in the same corpus and by the exact vacuum spectrum (Document~02).
\end{enumerate}

\section{Structural contributions worth preserving}

The following items appear to be genuine organizing ideas.
They are not theorems, but they are not nonsense.

\begin{enumerate}[leftmargin=*]
\item \textbf{Pipeline Architecture.}
A proposed chain:
\[
\text{(physical convexity/coercivity)} \Rightarrow \text{(LSI)} \Rightarrow \text{(spectral gap)} \Rightarrow \text{(RG stability)} \Rightarrow \text{(OS mass gap)}.
\]
This is a coherent architecture if each arrow is made precise with correct hypotheses (Document~03).

\item \textbf{Pipeline Architecture (computational).}
Use of a right-invariant tangent identification for $SU(3)^{E_L}$ together with an explicit gauge/toron projector to compute physically meaningful Hessian spectra, and (in principle) to transport Hessians through coarse-graining maps via Jacobians.

\item \textbf{Rigidity Mechanism (conjectural).}
The claim that a Riccati/Hamilton--Jacobi-type smoothing flow could ``convexify'' or stabilize the physical Hessian lower bound is a rigidity mechanism proposal.
It is not established; the curvature terms must be controlled.

\item \textbf{Implicit theorem (standard, but correctly recoverable).}
The vacuum quadratic form of Wilson action is the discrete Yang--Mills energy $\|dA\|^2$; the Hessian is a discrete $1$-form Laplacian (Document~02).
Standard result; appears in the lattice gauge theory literature.  No novelty claim supported.
\end{enumerate}

\section{Known vs new}

\begin{itemize}[leftmargin=*]
\item The Bakry--\'Emery log-Sobolev mechanism is a standard result (external).
\item The vacuum Hessian/Laplacian identification is standard.
\item The combination ``LSI + RG stability + OS'' as a single mass-gap proof pipeline:
\emph{I do not know} whether this exact packaging is present in the literature; the individual components are well known.
\item The HOTRG-Jacobian-based curvature propagation proposal:
plausibly novel as a computational program, but cannot be promoted without external verification and, in particular, without a correct treatment of differentiability/truncation and reflection positivity.
\end{itemize}

\section{What strengthens this work most}

\subsection*{Missing lemmas that would solidify the core}
\begin{enumerate}[leftmargin=*]
\item \textbf{Gauge reduction lemma.}
Construct a mathematically legitimate quotient or gauge-fixing framework producing a symmetric diffusion generator on the reduced space; show that it matches the intended ``physical'' projection.

\item \textbf{Non-vacuum, nonperturbative lower bound.}
Replace the (false-as-stated) uniform vacuum-neighborhood coercivity with a correct statement:
either a uniform convexity bound on the reduced space away from vacuum, or an effective-action lower bound after integrating UV modes, or a different coercivity notion that does not contradict Document~02.

\item \textbf{RG Lipschitz control (with truncation).}
Prove that the actual coarse-graining map used (e.g.\ HOTRG step with truncation) is Lipschitz on the relevant metric space and quantify its constant uniformly across scales, \emph{including} the SVD/truncation nonsmoothness.

\item \textbf{Reflection positivity preservation.}
Show that the coarse-graining step preserves (or at least does not destroy) the properties needed for OS reconstruction, and that the Euclidean spectral gap implied by the functional inequalities transfers to a Hamiltonian mass gap.
\end{enumerate}

\subsection*{Single sharpest theorem currently supported}
The strongest statement currently supported by the corpus is:

\medskip
\noindent
\emph{At $U\equiv I$, the Wilson action Hessian equals $(\beta/6)\Delta_1$ and its strictly physical finite-volume spectral gap equals $\frac{\beta}{6}\,4\sin^2(\pi/L)$.}
\medskip

\subsection*{Minimum additional assumption needed to close the main gap}
The minimal assumption that turns the pipeline into a valid theorem is:

\medskip
\noindent
\emph{Existence of a gauge-reduced diffusion generator on a smooth reduced space $Q_L$ such that the descended effective action is uniformly $\kappa$-convex with $\kappa>0$ independent of lattice size and scale.}
\medskip

This is strictly stronger than anything proved in the corpus, and it is not compatible with the vacuum linearization unless ``uniform'' refers to a different scaling regime or to a modified effective action.

\end{document}
